\section{$\Z/n\Z$: The Integers Modulo $n$}
$\mathbb{Z} =$ Set of integers.\\
Define a relation on $\mathbb{Z}$ by -
$$
a \sim b \text{ if } n \mid (a - b).
$$
Properties of this relation:
\begin{itemize}
    \item $a \sim a$ (Reflexivity)
    \item $a \sim b \implies b \sim a$ (Symmetry)
    \item $a \sim b, b \sim c \implies a \sim c$ (Transitivity)
\end{itemize}
Hence, this is an equivalence relation.\\
If $a \sim b$, we say $b$ is congruent to $a$ modulo $n$ and denote as $a \equiv b \text{ (mod }n)$.\\
Let $\overline{a} :=$ The set of all integers congruent to $a$ mod $n$, then $\overline{a} = \{a + kn: k \in \mathbb{Z}\}$ is called equivalence class / congruence class / residue class of $a$ mod $n$.\\
There are precisely $n$ distinct equivalence classes mod $n$, namely
$$\overline{0}, \overline{1}, \overline{2}, ..., \overline{n-1}.$$
\begin{eg}
    $\mathbb{Z}/4\mathbb{Z} = \{\overline{0}, \overline{1}, \overline{2}, \overline{3}\}.$
\end{eg}
Define an addition on $\mathbb{Z}/n\mathbb{Z}$ via -
$$\overline{a} \oplus \overline{b} = \overline{a+b}$$
\begin{eg}
    $\mathbb{Z}/3\mathbb{Z} = \{\overline{0}, \overline{1}, \overline{2}\}.$
    \[
    \begin{array}{c|ccc}
        \oplus       & \overline{0} & \overline{1} & \overline{2}\\
        \hline
        \overline{0} & \overline{0} & \overline{1} & \overline{2}\\
        \overline{1} & \overline{1} & \overline{2} & \overline{0}\\
        \overline{2} & \overline{2} & \overline{0} & \overline{1}\\
    \end{array}
    \]
\end{eg}
Notice that for all $a, b \in \mathbb{Z}/3\mathbb{Z}$ the following properties hold:
\begin{itemize}
    \item $(a \oplus b) \in \mathbb{Z}/3\mathbb{Z}$ (\textbf{Closure}).
    \item $(a \oplus b) \oplus c = a \oplus (b \oplus c)$ (\textbf{Associativity}).
    \item $a \oplus \overline{0} = \overline{0} \oplus a = a$ (\textbf{Identity element}).
    \item For all $a \in \mathbb{Z}/3\mathbb{Z}$, there exists $b \in \mathbb{Z}/3\mathbb{Z}$ such that $a \oplus b = b \oplus a = \overline{0}$ (\textbf{Inverse element}).
\end{itemize}
These properties hold for $\mathbb{Z}/n\mathbb{Z}$ in general.\\
Now let's move on to a totally different kind of set.
\section{Rotational Symmetries of an Equilateral Triangle}
Consider the set of all roational symmetries of an equilateral triangle.
\begin{figure}[H]
    \centering
    \incfig{symmetries-of-triangle}
    \caption{Symmetries of Triangle}
    \label{fig:symmetries-of-triangle}
\end{figure}
A symmetry is something we do to an object that leaves the object intact. In case of the equilateral triangle we have three roational symmetires: roatation by $0^{\circ}, 120^{\circ}$ and $240^{\circ}$. The operation here is to compose two symmetries, meaning apply one after another. Notice, composition of two symmetries is again a symmetry \textbf{(Closure)}. Composition of three symmetries is associative, because we really applying symmetries from the left to right \textbf{(Associativity)}. We consider rotation by $0^{\circ}$ or 'do nothing' is also a symmetry. When we compose  this with another symmetry, the other symmetry is unchanged \textbf{(Identity)}. Finally, given any symmetry, we can undo the symmetry \textbf{(Inverse elemnet)}.\\
Although, our earlier example of $\Z/3\Z$ and the Symmetries of a triangle is totally different set, both share four common properties.\\
A set equipped with an operation that follows these four properties is called a \textit{Group}.